\section{Relativization to a Class}

Relativization is essentially like satisfaction, just for classes. Our ultimate goal is the following: given $(A, R)$ a class structure and a formula $\psi$, we'd want to define $\psi^{(A, R)}$ to be the relativization (think ``interpretation'') of $\psi$ to $(A, R)$.

Here are some complications with this technique:
\begin{itemize}
    \item $R$ might not be $\in$
    \item There might be parameters involved in the definitions of $A$ and $R$. For instance, if $A$ is a set, then ``$A = \setst{x}{x \in A}$'' takes in a set $A$ as a parameter.
\end{itemize}

Look at $A = \setst{x}{\varphi(x)}$. As a scheme, define $\psi^A\of{\bar{v}}$ (with $\bar{v}$ denoting free variables) by recursion on the complexity of $\psi\of{\bar{v}}$ as follows. We will have atomic, propositional and quantifier steps.

Before starting, rename the variables of $\varphi$ other than $x$ so that none of them occurs in $\psi$. Write $\varphi'$ for the result, so that $A = \setst{x}{\varphi'(x)}$ still, and use $\varphi'$ in place of $\varphi$ throughout. Both directions of the clash matter. The quantifier step below substitutes a variable $u$ of $\psi$ for the free variable $x$, so a $u$ bound in $\varphi$ would be captured and $\varphi\of{u}$ would say something other than ``$u \in A$''; and what that substitution produces is then dropped underneath the quantifiers of $\psi$, so a parameter of $\varphi$ sharing its name with a variable that $\psi$ binds would be captured the other way.
\begin{itemize}
    \item \underline{Atomic:} We interpret equality and membership symbols as equality and membership. Ie,
    \begin{align*}
        \parenth{v_0 = v_1}^A &\text{ is } v_0 = v_1 \\
        \parenth{v_0 \in v_1}^A &\text{ is } v_0 \in v_1 
    \end{align*}

    \item \underline{Propositional:} If $\psi$ is a propositional combination of $\psi_0$ and $\psi_1$, then $\psi^A$ is the same propositional combination of $\psi_0^A$ and $\psi_1^A$.

    \item \underline{Quantifiers:} We relativize a quantifier by restricting it to $A$. Ie,
    \begin{align*}
        \parenth{\forall u \, \psi}^A &\text{ is } \forall u \brac{\varphi'(u) \to \psi^A} \\
        \parenth{\exists u \, \psi}^A &\text{ is } \exists u \brac{\varphi'(u) \land \psi^A}
    \end{align*}
\end{itemize}

The atomic step above is written for the case $R = \in$. For a general $R$ it is that step which changes, because ``$\in$'' is a symbol of the language we are interpreting, and what it is interpreted as in $(A, R)$ is $R$. So if $R = \setst{(x, y)}{\varphi_R\of{x, y, q}}$, then $\parenth{v_0 \in v_1}^{(A, R)}$ is $\varphi_R\of{v_0, v_1, q}$, and since the recursion bottoms out at the atomic subformulae, this replaces every occurrence of $\in$ in $\psi$. The parameter $q$ simply comes along for the ride, which is exactly why $\varphi_R$ needs the same renaming as $\varphi$ did, in both directions: in particular $q$ must not be a variable that $\psi$ binds.

Classes can be relativized as well as formulae, by relativizing the formula that defines them.

\begin{boxnotation}[Relativized Classes]
    Say $M$ is a transitive class and $A$ another class. Say $A = \setst{x}{\vp(x)}$. We might write $A^M$ to mean $\setst{x \in M}{\vp^M(x)}$ even though it may depend on the choice of $\varphi$ with this notation.
\end{boxnotation}
